\documentclass{chsmc}
\chsmcsetup{
  year = 2022,
  part = 1,
  type = B,
  show-answers = false,
}
\begin{document}

\section{填空题：本大题共 8 小题，每小题 8 分，共 64 分.}

\begin{enumerate}
  \item 不等式 $\dfrac{20}{x-9} > \dfrac{22}{x-11}$ 的解集为 \blankline
  \item 在平面直角坐标系中，以抛物线 $\Gamma\colon y^2 = 6x$ 的焦点为圆心作一个圆
    $\Omega$，与 $\Gamma$ 的准线相切，则圆 $\Omega$ 的面积为 \blankline
  \item 函数 $f(x) = \lg 2\cdot\lg 5 - \lg 2x\cdot\lg 5x$ 的最大值为 \blankline
  \item 一枚不均匀的硬币，若随机抛掷它两次均得到正面的概率为 $\dfrac{1}{2}$，
    则随机抛掷它两次得到正面、反面各一次的概率为 \blankline
  \item 已知复数 $z$ 满足 $|z| = 1$，且
    $\operatorname{Re}\dfrac{z+1}{\bar{z}+1} = \dfrac{1}{3}$，
    则 $\operatorname{Re}\dfrac{z}{\bar{z}}$ 的值为 \blankline
  \item 若正四棱锥 $P-ABCD$ 的各条棱长均相等，$M$ 为棱 $AB$ 的中点，
    则异面直线 $BP$ 与 $CM$ 所成的角的余弦值为 \blankline
  \item 若 $\triangle ABC$ 的三个内角 $A, B, C$ 满足
    $\cos A = \sin B = 2\tan\dfrac{C}{2}$，
    则 $\sin A + \cos A + 2\tan A$ 的值为 \blankline
  \item 一个单位方格的四条边中，若存在三条边染了三种不同的颜色，
    则称该单位方格是“多彩”的. 如图，一个 $1 \times 3$ 方格表的表格线共含
    $10$ 条单位长线段，现要对这 $10$ 条线段染色，每条线段染为红、黄、蓝三色之一，
    使得三个单位方格都是多彩的. 这样的染色方式数为 \blankline （答案用数值表示）
\end{enumerate}

\section{解答题：本大题共 3 个小题，满分 56 分. 解答应写出文字说明、证明过程或演算步骤.}

\begin{enumerate}[resume]
  \item (本题满分 16 分) 在平面直角坐标系中，$F_1$, $F_2$ 是双曲线
    $\Gamma\colon \dfrac{x^2}{3} - y^2 = 1$ 的两个焦点，$\Gamma$ 上一点 $P$
    满足 $\overrightarrow{PF_1}\cdot\overrightarrow{PF_2} = 1$.
    求点 $P$ 到 $\Gamma$ 的两条渐近线的距离之和.
  \item (本题满分 20 分) 设正数 $a_1, a_2, a_3, b_1, b_2, b_3$ 满足：
    $a_1, a_2, a_3$ 成公差为 $b_1$ 的等差数列，$b_1, b_2, b_3$ 成公比为 $a_1$
    的等比数列，且 $a_3 = b_3$. 求 $a_3$ 的最小值，并确定当 $a_3$ 取到最小值时
    $a_2 b_2$ 的值.
  \item (本题满分 20 分) 若 $a$, $b$ 为实数，$a < b$，函数 $y = \sin x$
    在闭区间 $[a,b]$ 上的最大值与最小值之差为 $1$，求 $b - a$ 的取值范围.
\end{enumerate}

\end{document}